\section{Methodology}
\label{sec:methodology}

\subsection{The DML--DID Framework}

We employ a partially linear regression (PLR) model in the DML
framework of \citet{chernozhukov2022automatic}. The model
specifies the outcome $Y$ as a function of the treatment $D$ and a
set of confounders $X$:
\begin{equation}
\label{eq:plr}
Y = \theta \cdot D + g(X) + \varepsilon, \qquad
D = m(X) + \nu,
\end{equation}
where $g(\cdot)$ and $m(\cdot)$ are unknown nuisance functions,
$\theta$ is the causal parameter of interest, and
$\mathbb{E}[\varepsilon \mid X, D] = 0$,
$\mathbb{E}[\nu \mid X] = 0$.

The treatment variable $D$ is constructed as the interaction between
a policy dummy ($\did_i$) and a continuous climate shock intensity
measure ($S_{it}$):
\begin{equation}
\label{eq:treatment}
D_{it} = \did_{it} \times S_{it},
\end{equation}
where $\did_{it} = \mathbf{1}[t \geq 2017] \times \mathbf{1}[i \in
\mathcal{T}]$ and $\mathcal{T}$ denotes the set of 27 pilot cities.
The shock intensity $S_{it}$ is standardised to $[0, 1]$ and captures
the severity of extreme weather events experienced by city $i$ in
year $t$, enabling identification of the policy's differential effect
under climate stress.

\subsection{Cross-Fitting and Nuisance Estimation}

To estimate $\theta$ without parametric restrictions on $g$ and $m$,
we implement $K$-fold cross-fitting. The sample is partitioned into
$K$ mutually exclusive folds. For each fold $k$, nuisance functions
are estimated on the remaining $K-1$ folds and used to predict
residualised outcomes and treatments on the held-out fold:
\begin{align}
\tilde{Y}_i &= Y_i - \hat{g}_k(X_i), \\
\tilde{D}_i &= D_i - \hat{m}_k(X_i).
\end{align}
The final estimator is the residual-on-residual OLS regression:
\begin{equation}
\label{eq:dml_estimator}
\hat{\theta} = \frac{\sum_i \tilde{D}_i \tilde{Y}_i}
                    {\sum_i \tilde{D}_i^2}.
\end{equation}

Before cross-fitting, all variables are within-transformed by
demean-ing over city and year fixed effects to absorb time-invariant
unobservables and common temporal shocks, following the TWFE--DML
approach of \citet{ahrens2024ddml}. The nuisance functions are
estimated using four machine-learning learners---Random Forest (RF),
XGBoost, Neural Network (NN), and LASSO---enabling comparison across
algorithms as a specification robustness check.

\subsection{Inference}

Standard errors are computed using the heteroskedasticity-robust
DML variance formula \citep{chernozhukov2022automatic}:
\begin{equation}
\label{eq:se}
\widehat{\mathrm{SE}}(\hat{\theta}) =
\frac{\sqrt{\sum_i (\tilde{D}_i \hat{\varepsilon}_i)^2}}
     {\sum_i \tilde{D}_i^2},
\end{equation}
where $\hat{\varepsilon}_i = \tilde{Y}_i - \hat{\theta}\tilde{D}_i$.
For the main results table, we additionally report bootstrap standard
errors based on 200 resampling iterations of the residualised data.

\subsection{CSEE Construction}
\label{sec:csee}

The Climate-Stress Ecological Elasticity (CSEE) index decomposes
ecological resilience into two components:
\begin{equation}
\CSEE_{it} = 0.5 \times \mathrm{CR}_{it} + 0.5 \times \mathrm{RC}_{it},
\end{equation}
where CR (Climate Resistance) measures the degree to which NDVI is
maintained during extreme weather events, and RC (Recovery) measures
the speed of NDVI restoration afterwards.

\paragraph{Climate Resistance (CR).}
For each extreme weather event $e$ affecting city $i$, we compute:
\begin{equation}
\mathrm{CR}_{ie} = 1 - \frac{|\mathrm{NDVI}_{ie}^{\mathrm{event}} -
       \mathrm{NDVI}_{i}^{\mathrm{normal}}|}
      {\mathrm{NDVI}_{i}^{\mathrm{normal}}},
\end{equation}
where $\mathrm{NDVI}_{i}^{\mathrm{normal}}$ is the 5-year pre-event
mean, and $\mathrm{NDVI}_{ie}^{\mathrm{event}}$ is the NDVI during the
event window (30 days centred on the event). CR is then averaged across
all events in city-year $it$.

\paragraph{Climate Recovery (RC).}
Recovery is measured as the fraction of NDVI loss regained within a
3-month post-event window:
\begin{equation}
\mathrm{RC}_{ie} = \frac{\mathrm{NDVI}_{ie}^{\mathrm{post}} -
       \mathrm{NDVI}_{ie}^{\mathrm{event}}}
      {\mathrm{NDVI}_{i}^{\mathrm{normal}} -
       \mathrm{NDVI}_{ie}^{\mathrm{event}}},
\end{equation}
where $\mathrm{NDVI}_{ie}^{\mathrm{post}}$ is the mean NDVI in the
3 months following the event. Cross-year recovery is attributed to the
year in which the event began, accommodating events that span calendar
year boundaries.

\subsection{Parallel Trends and Event Study}

Identification in the DID--DML framework rests on the parallel-trends
assumption: in the absence of treatment, outcome trends in pilot and
control cities would have evolved similarly. We test this assumption
using an event-study specification with formal Wald $F$-tests on
pre-treatment leads:
\begin{equation}
\label{eq:event_study}
Y_{it} = \alpha_i + \delta_t + \sum_{k=-6, k\neq -1}^{6}
         \beta_k \cdot \mathbf{1}[t - 2017 = k] \cdot
         \mathbf{1}[i \in \mathcal{T}] + \gamma X_{it} + \varepsilon_{it},
\end{equation}
where $k = -1$ is the reference period. The null hypothesis
$H_0: \beta_{-6} = \cdots = \beta_{-2} = 0$ is tested via a joint
Wald $F$-statistic, and failure to reject at the 10\% level is
interpreted as support for parallel trends
\citep{rambachan2023more}.

\subsection{Robustness Design}

We implement six robustness checks. (i)~Alternative ML algorithms
(RF, XGBoost, NN, LASSO) to assess sensitivity to learner choice.
(ii)~A 200-iteration placebo test, randomly permuting treatment
assignment within each year and estimating $\hat{\theta}$ on each
permutation; the placebo $p$-value is the fraction of permutations
yielding $|\hat{\theta}_{\mathrm{placebo}}| \geq
|\hat{\theta}_{\mathrm{true}}|$. (iii)~Propensity-score matching
(PSM)--DML, matching treated cities to controls on pre-treatment
covariates before DML estimation. (iv)~Subsample restrictions,
excluding COVID-affected years (2020--2022) and restricting to
post-2010 and post-2008 samples. (v)~Cross-fitting sensitivity with
$K \in \{3, 5, 7, 10\}$. (vi)~Lagged treatment effects up to 3 years
to capture delayed policy impacts.
